(1) 单个像素的复电场由 \(N\ge2\) 个幅值相等、统计独立且在 \([0,2\pi)\) 均匀分布的相位波元组成，
\[
E_N=\frac{A}{\sqrt N}\sum_{j=1}^N e^{i\phi_j},\quad I_N=|E_N|^2.
\]
不预先假设高斯近似，从随机相位的精确二阶与四阶矩求平均强度、方差及有限 \(N\) 的散斑对比度 \(C_N\)。
(2) 当 \(N\to\infty\) 时根据复中心极限定理求圆对称复高斯电场的强度概率密度与累积分布，检验 \(\langle I^2\rangle=2\langle I\rangle^2\) 和对比度极限。
(3) 对两像素零均值联合圆复高斯场 \(E_1,E_2\)，定义互相干度 \(g^{(1)}_{12}=\langle E_1^*E_2\rangle/\sqrt{\langle I_1\rangle\langle I_2\rangle}\)。使用复高斯四阶配对独立证明归一化强度关联 \(g^{(2)}_{12}=1+|g^{(1)}_{12}|^2\)，检查完全相关与不相关边界。
(4) 若探测器一次曝光平均了 \(M\) 个等统计、相互独立的充分发展散斑强度，求平均强度的精确Gamma概率密度及对比度 \(C_M\)，说明何时曝光时间内不独立或存在相干背景会破坏此结论；区分经典散斑涨落与光子泊松噪声。
